\documentclass[10pt,a4paper]{amsart}
\usepackage[margin=19mm]{geometry}
\usepackage{amsmath,amssymb,mathtools}
\usepackage[scr=rsfs,cal=euler]{mathalfa}
\usepackage{xfrac}
\usepackage[unicode,hidelinks,bookmarks=false]{hyperref}
\newcommand{\ie}{i.e.\ }
\newcommand{\cf}{cf.\ }
\newcommand{\resp}{respectively\ }
\newcommand{\Subsection}{\subsection}
\providecommand{\nobreakdash}{\nobreak\hbox{-}}
\setlength{\emergencystretch}{2em}
\multlinegap0pt
\usepackage{bm}
\usepackage{bbm}
\usepackage{dsfont}
\usepackage{xspace}
\usepackage{cancel}
\usepackage{mathtools}
\usepackage{amsthm}
\makeatletter
\@ifundefined{prop}{\newtheorem{prop}{Prop}}{}
\makeatother
\newcommand{\der}{\textup{d}}
\title[Mass concentration in a spatially inhomogeneous coagulation ]{Mass concentration in a spatially inhomogeneous coagulation model with fast sedimentation}
\author{Iulia Cristian, Juan J. L. Vel\'azquez}
\date{Source: arXiv:2510.04270v1 (CC BY 4.0)}
\begin{document}
\maketitle

\noindent\emph{Source and licence.} Mass concentration in a spatially inhomogeneous coagulation model with fast sedimentation. Version arXiv:2510.04270v1 (CC BY 4.0), distributed under the Creative Commons Attribution 4.0 International licence, as recorded in the arXiv licence metadata for this exact version and shown on the abstract page https://arxiv.org/abs/2510.04270v1. Full rights notice: licenses/arxiv-2510.04270-CC-BY-4.0.txt.\par\smallskip

\noindent\textit{Editorial scope.} Counted unit: the Prop whose source label is \texttt{der v g negative} in the article (source0.tex lines 2296--2301, source label \texttt{der v g negative}), delivered together with the article's own complete original proof of it (source0.tex lines 2302--2478, one \texttt{proof} environment of 177 lines). No mathematics is written, repaired or re-derived: statement and proof are the article's own text line for line. Required context retained uncounted: initial condition decay (lines 495--497); equation supersolution (lines 2058--2060); we need function to be even (lines 2141--2146); the original characteristics (lines 2148--2153); proposition about inequalities (lines 2160--2168); ineq (lines 2162--2167); iterative argument characteristics (lines 2190--2192). Every definition, display and label of those lines is retained unchanged. Omitted from this package and not redistributed: 1--494, 498--2057, 2061--2140, 2147--2147, 2154--2159, 2168--2189, 2193--2295, 2479--3293. Nothing inside a retained excerpt is omitted or altered: every retained body is delivered exactly as the article's lines are written, so no omission receipt of any kind accompanies this package. Citations: the retained excerpts carry no citation command at all, so no bibliography entry of the article is transcribed and no reference is redistributed. Reference adaptation: none was needed and none was made; every reference inside the delivered unit resolves inside it, so no unresolved reference or citation is produced. Printed slips kept literal: no printed word, symbol, hypothesis or number of the counted statement or of its proof was altered. Layout and class: the article's own class is \texttt{article} with packages including inputenc, amscd, amsmath, amsthm, amssymb, amsfonts, txfonts, eucal, bbm, appendix, authblk, graphicx, pifont, hyperref; the wrapper is the lane's pinned \texttt{amsart} editorial wrapper at 10pt on A4, which restates the theorem-like environments the retained text needs and adds the article's own macro definitions that the retained text needs (\texttt{\textbackslash der}), with extra wrapper packages (bm, bbm, dsfont, xspace, cancel, mathtools). The delivered numbering is the unit's own. Assets: no retained span contains a figure environment, a graphics-inclusion command, a PDF-page inclusion, a table or a TikZ picture; no image file is copied into this package, the package declares no asset dependency and no image file is redistributed with it. Licence: version arXiv:2510.04270v1 of this article is distributed under the Creative Commons Attribution 4.0 International licence, as recorded in the arXiv licence metadata for this exact version and shown on the abstract page https://arxiv.org/abs/2510.04270v1. A full rights notice is delivered at licenses/arxiv-2510.04270-CC-BY-4.0.txt. Characters: every retained excerpt is pure ASCII, so the delivered engine, XeLaTeX with its default Latin Modern text font, reports no missing character.
\begin{align}\label{initial condition decay}
   (1+v^{b}) H_{\epsilon,\textup{in}}(y,v)\leq \frac{A}{1+|y-v^{\alpha}|^{m}},
\end{align}

\begin{align}\label{equation supersolution}
      \partial_{t}G_{\epsilon}(y,v,t)+\frac{1}{\epsilon}(v^{\alpha}-y)\partial_{y}G_{\epsilon}(y,v,t)+\frac{Lv^{\gamma}\xi(v)}{1+|y|^{d}}\partial_{v}G_{\epsilon}(y,v,t)-\frac{1}{\epsilon}G_{\epsilon}(y,v,t)=0.
\end{align}

  \begin{equation}\label{we need function to be even}
\left\{\begin{aligned}
d&=\bigg\lfloor\frac{b-2}{\alpha}\bigg\rfloor-1; &\textup{ if }\bigg\lfloor\frac{b-2}{\alpha}\bigg\rfloor& \textup{ odd; }\\
d&=\bigg\lfloor\frac{b-2}{\alpha}\bigg\rfloor;&\textup{ if }\bigg\lfloor\frac{b-2}{\alpha}\bigg\rfloor& \textup{ even, }
   \end{aligned}\right.
   \end{equation}

  \begin{equation}\label{the original characteristics}
\left\{\begin{aligned}
\partial_{t}{Y_{\epsilon}}(y,v,t)&=-\frac{1}{\epsilon}({V_{\epsilon}}^{\alpha}-Y), & \qquad {Y_{\epsilon}}(y,v,0)&=y\leq \big(\frac{v}{3}\big)^{\alpha}\,, \\
\partial_{t}{V_{\epsilon}}(y,v,t)&=-\frac{{L}{V_{\epsilon}}^{\gamma}\xi(V_{\epsilon})}{1+|{Y_{\epsilon}}|^{d}},& \qquad {V_{\epsilon}}(y,v,0)&=v.
   \end{aligned}\right.
   \end{equation}

   \begin{prop}\label{proposition about inequalities}
Let $\alpha\in(0,1)$, $\gamma\in(1,1+\alpha)$, and $d>1$ even. Let $L>0$ and let $Y_{\epsilon}, V_{\epsilon}$ be as in \eqref{the original characteristics}. Then there exists a sufficiently small $\epsilon_{1}\in(0,1)$, which depends on $L$, such that for all $t \geq 0$ and $\epsilon\leq \epsilon_{1}$, the following estimates hold for all $v>0$ and $y\leq \big(\frac{v}{3}\big)^{\alpha}$.
      \begin{align}
        \frac{9}{10} v &\leq  V_{\epsilon}(y,v,t)\leq v;\label{ineq}\\
     Y_{\epsilon}(y,v,t)&\leq \big(\frac{v}{3}\big)^{\alpha};\label{ineq part two}\\
e^{\frac{t}{\epsilon}} (y-v^{\alpha}) & \leq   Y_{\epsilon}-V_{\epsilon}^{\alpha};\label{ineqq}\\
 Y_{\epsilon}-V_{\epsilon}^{\alpha}&\leq e^{\frac{t}{\epsilon}} (y-\Big(\frac{9}{10}\Big)^{\alpha}v^{\alpha}).\label{ineq 3}
      \end{align}
  \end{prop}

  \begin{align}\label{iterative argument characteristics}
  \partial_{s}{Y}(y,v,s)&=\frac{1}{\epsilon}(Y-{V}^{\alpha})\leq \frac{1}{\epsilon}(Y-\Big(\frac{9}{10}\Big)^{\alpha}{v}^{\alpha}).
  \end{align}

 \begin{prop}\label{der v g negative}
Fix $\alpha\in(0,1)$, $\gamma\in(1,1+\alpha)$, $L>0$ and $\epsilon\in(0,1)$. Let $d>1$ in \eqref{we need function to be even} even and sufficiently large such that $\gamma\leq \frac{d\alpha}{d+1}+1$. Let $G_{\epsilon,L}$ solve equation \eqref{equation supersolution}. Then there exists a sufficiently small $\epsilon_{1}\in(0,1)$ such that for all $t \geq 0$ and $\epsilon\leq \epsilon_{1}$, we have that for all $v>0$ and $y\leq \big(\frac{v}{3}\big)^{\alpha}$,
\begin{align*}
    \partial_{v}G_{\epsilon,L}(y,v,t)\leq 0.
\end{align*}
 \end{prop}

 \begin{proof}
  Let $Y_{\epsilon}, V_{\epsilon}$ be as in \eqref{the original characteristics}.    We have that 
     \begin{align*}
         G_{\epsilon,L}(y,v,t)=\frac{A e^{\frac{t}{\epsilon}}}{(1+V_{\epsilon}^{b})(1+|Y_{\epsilon}-V_{\epsilon}^{\alpha}|^{m})},
     \end{align*}
with $A$ as in \eqref{initial condition decay}. As before, in order to simplify the notation, we replace the notation of the pair $(Y_{\epsilon},V_{\epsilon})$ in \eqref{the original characteristics} with $(Y,V)$, while keeping in mind the dependence on $\epsilon$.    Moreover, from Proposition \ref{proposition about inequalities}, we know that $Y-V^{\alpha}\leq 0$, for all $t\geq 0$. Thus
     \begin{align*}
         \partial_{v}G_{\epsilon,L}(y,v,t)=-\frac{mA e^{\frac{t}{\epsilon}}(V^{\alpha}-Y)^{m-1}[\alpha V^{\alpha-1}\partial_{v}V-\partial_{v}Y]}{(1+V^{b})(1+(V^{\alpha}-Y)^{m})^{2}}-\frac{b Ae^{\frac{t}{\epsilon}}V^{b-1}\partial_{v}V}{(1+V^{b})^{2}(1+(V^{\alpha}-Y)^{m})}.
     \end{align*}
     Thus, in order to prove that $\partial_{v}G_{\epsilon,L}(y,v,t)\leq 0$ it suffices to prove that $\alpha V^{\alpha-1}\partial_{v}V-\partial_{v}Y\geq 0$ and $\partial_{v}V\geq 0$. Fix $t>0$. We know that $V\geq 0$ by Proposition \ref{proposition about inequalities}. We will prove that
     \begin{align*}
         \partial_{v}V(y,v,t)\geq 0, \textup{ for all } v>0, y\leq (\frac{v}{3})^{\alpha},
     \end{align*}
     and that
     \begin{align*}
         \partial_{v}Y(y,v,t)\leq 0, \textup{ for all } v>0, y\leq (\frac{v}{3})^{\alpha},
     \end{align*}
     thus concluding our proof. Since $\partial_{v}V(y,v,0)=1$, by continuity there exists $t_{1}>0$ such that $\partial_{v}V(y,v,s)\in[\frac{1}{2},1]$, for all $s\in[0,t_{1}]$. By \eqref{the original characteristics}, we have that
     \begin{align*}
         \partial_{s}\partial_{v}Y(y,v,s)=\frac{\partial_{v}Y}{\epsilon}-\frac{\alpha V^{\alpha-1}\partial_{v}V}{\epsilon}, \qquad \partial_{v}Y(y,v,0)=0.
     \end{align*}
     In other words, using Proposition \ref{proposition about inequalities} and since $\alpha<1$, it holds that for all $s\in[0,t_{1}]$ 
     \begin{align*}
             \partial_{s}\big(e^{-\frac{s}{\epsilon}}\partial_{v}Y(y,v,s)\big)=-e^{-\frac{s}{\epsilon}}\frac{\alpha V^{\alpha-1}\partial_{v}V}{\epsilon}\leq -e^{-\frac{s}{\epsilon}}\frac{\alpha v^{\alpha-1}}{2\epsilon}.
     \end{align*}
     Thus, by integrating in time, we obtain that
     \begin{align*}
         e^{-\frac{s}{\epsilon}}\partial_{v}Y\leq \frac{\alpha}{2} e^{-\frac{s}{\epsilon}}v^{\alpha-1}-\frac{\alpha}{2} v^{\alpha-1},
     \end{align*}
     for all $s\in[0,t_{1}],$ and thus
     \begin{align*}
         \partial_{v}Y(y,v,s)\leq \frac{\alpha}{2} v^{\alpha-1}(1-e^{\frac{s}{\epsilon}})\leq 0.
     \end{align*}
Using similar arguments, we obtain that
\begin{align}\label{exponential bound for der y}
    \partial_{v}Y(y,v,s)\geq \frac{10^{1-\alpha}\alpha}{9^{1-\alpha}} v^{\alpha-1}(1-e^{\frac{s}{\epsilon}}).
\end{align}
%By assumption, $y\not\in[\big((\frac{9}{10})^{\alpha}-\delta\big)v^{\alpha}(1-e^{-\frac{t}{\epsilon}}),(1+\delta)v^{\alpha}(1-e^{-\frac{t}{\epsilon}})]$. If $y\leq \big((\frac{9}{10})^{\alpha}-\delta\big)v^{\alpha}(1-e^{-\frac{t}{\epsilon}})$ and $s<t$, then $y\leq  \big((\frac{9}{10})^{\alpha}-\delta\big)v^{\alpha}(1-e^{-\frac{s}{\epsilon}})$.  By Proposition \ref{proposition about inequalities}, we have that
%\begin{align*}
   % Y(y,v,s)\leq e^{\frac{s}{\epsilon}}\big(y-(\frac{9}{10})^{\alpha}v^{\alpha}\big)+(\frac{9}{10})^{\alpha}v^{\alpha}&\leq  e^{\frac{s}{\epsilon}}\big[((\frac{9}{10})^{\alpha}-\delta)v^{\alpha}(1-e^{-\frac{s}{\epsilon}})-(\frac{9}{10})^{\alpha}v^{\alpha}(1-e^{-\frac{s}{\epsilon}})\big]\\
   % &=-\delta v^{\alpha}(e^{\frac{s}{\epsilon}}-1).
%\end{align*}
%If $y\geq (1+\delta)v^{\alpha}(1-e^{-\frac{t}{\epsilon}})$, since $y\leq (\frac{v}{2})^{\alpha}$ ans $s<t$, it holds that
%\begin{align*}
  %  Y(y,v,s)\geq e^{\frac{s}{\epsilon}}\big(y-v^{\alpha}\big)+v^{\alpha}\geq e^{\frac{t}{\epsilon}}\big(y-v^{\alpha}\big)+v^{\alpha}&\geq e^{\frac{t}{\epsilon}}\big[(1+\delta)v^{\alpha}(1-e^{-\frac{t}{\epsilon}})-v^{\alpha}(1-e^{-t}{\epsilon}\big)\big]\\
 %   &= \delta v^{\alpha}(e^{\frac{t}{\epsilon}}-1)\geq  v^{\alpha}(e^{\frac{s}{\epsilon}}-1).
%\end{align*}
%Thus, in both cases we have that
%\begin{align}\label{absolute value y bound}
 %   |Y(y,v,s)|\geq \delta v^{\alpha}(e^{\frac{s}{\epsilon}}-1).
%\end{align}

By \eqref{the original characteristics}, we also have that for all $s\in[0,t_{1}]$,
\begin{align}\label{der v v}
    \partial_{s}\partial_{v}V(y,v,s)=-\frac{\gamma L V^{\gamma-1}\partial_{v}V}{1+|Y|^{d}}+\frac{d L V^{\gamma}|Y|^{d-2}Y\partial_{v}Y}{(1+|Y|^{d})^{2}}, \qquad \partial_{v}V(y,v,0)=1.
\end{align}
We will prove that 
\begin{align}\label{linear ode part 1}
    \int_{0}^{s}\frac{ V^{\gamma-1}(z)\der z}{1+|Y|^{d}(z)}\leq C\epsilon
\end{align}
and that 
\begin{align}\label{linear ode part 2}
    \int_{0}^{s}\frac{V^{\gamma}|Y|^{d-1}|\partial_{v}Y|}{(1+|Y|^{d})^{2}}\der z \leq C\epsilon,
\end{align}
for all $s\in[0,t_{1}]$. We first prove that \eqref{linear ode part 1} holds. By \eqref{ineq} we have that
\begin{align*}
   \frac{V^{\gamma-1}}{1+|Y|^{d}}\leq \frac{ v^{\gamma-1}}{1+|Y|^{d}}.
\end{align*}
Moreover, we remember that 
\begin{align*}
    \partial_{s}Y(y,v,s)=\frac{Y}{\epsilon}-\frac{V^{\alpha}}{\epsilon}\leq \frac{v^{\alpha}}{\epsilon}\bigg[\frac{1}{3^{\alpha}}-\Big(\frac{9}{10}\Big)^{\alpha}\bigg]\leq -\frac{c v^{\alpha}}{\epsilon}
\end{align*}
and then as before we have the following cases. \begin{enumerate}\item[i)]  $Y(y,v,0)=y\in[0,(\frac{v}{3})^{\alpha}]$ and $Y(y,v,t_{1})< 0$.
\end{enumerate}
Then there exists $\tau:=\tau(y,v)\in[0,t_{1})$ such that $Y(y,v,\tau(y,v))=0$ and, for all $s\in[0,t_{1}]$, it holds that
\begin{align}\label{linear bound for y}
 |Y(y,v,s)|\geq \frac{c v^{\alpha}}{\epsilon}|s- \tau(y,v)|.
\end{align}
Then 
\begin{align*}
        \int_{0}^{s}\frac{ V^{\gamma-1}(z)\der z}{1+|Y|^{d}(z)}\leq      \int_{0}^{s}\frac{ Cv^{\gamma-1}\der z}{1+ \frac{ v^{\alpha d}}{\epsilon^{d}}|z- \tau(y,v)|^{d}}.
\end{align*}
We make the change of variable $\xi=\frac{v^{\alpha}}{\epsilon}(z-\tau)$ and use the fact that we can consider without loss of generality that $v\geq 1$ in order to obtain that
\begin{align*}
        \int_{0}^{s}\frac{ V^{\gamma-1}(z)\der z}{1+|Y|^{d}(z)}\leq      \int_{-\infty}^{\infty}\frac{ C\epsilon v^{\gamma-\alpha-1}\der \xi}{1+ |\xi|^{d}}\leq C \epsilon.
\end{align*}
\begin{enumerate}
    \item[ii)]If $Y(y,v,0)=y\in[0,(\frac{v}{3})^{\alpha}]$ and $Y(y,v,t_{1})\geq 0$, we have that \begin{align*}
    |Y(y,v,s)|=Y(y,v,s)\geq Y(y,v,s)-Y(y,v,t_{1})=|Y(y,v,s)-Y(y,v,t_{1})|\geq \frac{c v^{\alpha}}{\epsilon}|s- t_{1}|
\end{align*}\end{enumerate}and the result follows as before. 

\begin{enumerate}\item[iii)]Otherwise $Y(y,v,0)=y<0$. Then for all $s\in[0,t_{1}]$, it holds by  \eqref{iterative argument characteristics} that
  \begin{align*}
Y(y,v,s)\leq -\Big(\frac{9}{10}\Big)^{\alpha}\frac{v^{\alpha}s}{\epsilon}
  \end{align*}\end{enumerate}and proceed as in the previous case in order to show that \eqref{linear ode part 1} holds.


We now prove \eqref{linear ode part 2}. We only analyze the case when $Y(y,v,0)=y\in[0,(\frac{v}{3})^{\alpha}]$ and $Y(y,v,t_{1})<0$ as the other two cases can be treated in a similar, more straightforward manner.
Thus, let us assume $Y(y,v,0)=y\in[0,(\frac{v}{3})^{\alpha}]$ and $Y(y,v,t_{1})<0$. Then, for every $y$ and every $v>0$, there exists $\tau:=\tau(y,v)\in[0,t_{1})$ such that $Y(y,v,\tau(y,v))=0$. We can write 
\begin{align*}
     \int_{0}^{s}\frac{V^{\gamma}|Y|^{d-1}|\partial_{v}Y|}{(1+|Y|^{d})^{2}}\der z = \int_{0}^{\tau(y,v)}\frac{V^{\gamma}|Y|^{d-1}|\partial_{v}Y|}{(1+|Y|^{d})^{2}}\der z + \int_{\tau(y,v)}^{s}\frac{V^{\gamma}|Y|^{d-1}|\partial_{v}Y|}{(1+|Y|^{d})^{2}}\der z =:(I)+(II).
\end{align*}
We bound $(II)$.
We have that
\begin{align*}
    \partial_{s}Y=\frac{Y}{\epsilon}-\frac{V^{\alpha}}{\epsilon}\leq \frac{Y}{\epsilon}-(\frac{9}{10})^{\alpha}\frac{ v^{\alpha}}{\epsilon}
\end{align*}
and thus
\begin{align*}
    \partial_{s}(e^{-\frac{s}{\epsilon}}Y)\leq -e^{-\frac{s}{\epsilon}}(\frac{9}{10})^{\alpha}\frac{v^{\alpha}}{\epsilon}.
\end{align*}
We integrate from $\tau(y,v)$ to $s\in[\tau(y,v),t_{1}]$ and use the fact that $Y(y,v,\tau(y,v))=0$ in order to obtain that
\begin{align*}
   e^{-\frac{s}{\epsilon}}Y(s)\leq e^{-\frac{s}{\epsilon}}(\frac{9}{10})^{\alpha}v^{\alpha}-e^{-\frac{\tau}{\epsilon}}(\frac{9}{10})^{\alpha}v^{\alpha}
\end{align*}
or in other words
\begin{align}\label{exponential bound for y}
    Y(y,v,s)\leq (\frac{9}{10})^{\alpha}v^{\alpha}(1-e^{\frac{s-\tau}{\epsilon}})\leq 0.
\end{align}
Combining \eqref{exponential bound for y} with \eqref{exponential bound for der y}, we have that
\begin{align*}
 \bigg|\frac{V^{\gamma}|Y|^{d-1}\partial_{v}Y}{(1+|Y|^{d})^{2}}\bigg|&\leq  \frac{C  v^{\gamma}|\partial_{v}Y|}{1+|Y|^{d+1}}\leq  \frac{C  v^{\gamma}|\partial_{v}Y|}{1+C v^{\alpha (d+1)}(e^{\frac{s-\tau}{\epsilon}}-1)^{d+1}}\\
 &\leq \frac{C  v^{\gamma+\alpha-1}(e^{\frac{s}{\epsilon}}-1)}{1+v^{\alpha (d+1)}(e^{\frac{s-\tau}{\epsilon}}-1)^{d+1}}\\
 &\leq \frac{C  v^{\gamma-1} (e^{\frac{s}{\epsilon}}-1)}{1+ v^{\alpha d}(e^{\frac{s-\tau}{\epsilon}}-1)^{d+1}}.
\end{align*}
Thus
\begin{align*}
\int_{\tau}^{s} \bigg|\frac{V^{\gamma}|Y|^{d-2}Y\partial_{v}Y}{(1+|Y|^{d})^{2}}\bigg|\der z \leq \int_{\tau}^{s}\frac{C  v^{\gamma-1} (e^{\frac{z}{\epsilon}}-1)}{1+ v^{\alpha d}(e^{\frac{z-\tau}{\epsilon}}-1)^{d+1}}\der z.
\end{align*}
We make the change of variables $\xi=v^{\frac{d\alpha }{d+1}}(e^{\frac{z-\tau}{\epsilon}}-1)$ and thus
\begin{align*}
    \der \xi = \frac{v^{\frac{\alpha d}{d+1}}}{\epsilon}e^{\frac{z-\tau}{\epsilon}}\der z.
\end{align*}
We have that  $\tau\leq C\epsilon$. This is since if we choose $s=0$ in \eqref{linear bound for y}, then $|y|\geq \frac{c v^{\alpha}}{\epsilon}| \tau(y,v)|$ and we are in the case when $y\in[0, (\frac{v}{3})^{\alpha}]$. Using that $\tau\leq C\epsilon$ and $v\geq 1$, we further obtain that
\begin{align}\label{bound ii}
\int_{\tau}^{s} \bigg|\frac{V^{\gamma}|Y|^{d-2}Y\partial_{v}Y}{(1+|Y|^{d})^{2}}\bigg|\der z &\leq \int_{0}^{\infty}\frac{C \epsilon  v^{\gamma-\frac{\alpha d}{d+1}-1}e^{-\frac{z-\tau}{\epsilon}}(e^{\frac{z}{\epsilon}}-1) }{1+ |\xi|^{d+1}}\der \xi\leq \int_{0}^{\infty}\frac{C \epsilon  v^{\gamma-\frac{\alpha d}{d+1}-1}e^{\frac{\tau}{\epsilon}}(1-e^{-\frac{z}{\epsilon}}) }{1+ |\xi|^{d+1}}\der \xi\nonumber\\
&\leq \int_{0}^{\infty}\frac{C \epsilon  v^{\gamma-\frac{\alpha d}{d+1}-1}e^{\frac{\tau}{\epsilon}}}{1+ |\xi|^{d+1}}\der \xi\leq\int_{0}^{\infty}\frac{C \epsilon  v^{\gamma-\frac{\alpha d}{d+1}-1}}{1+ |\xi|^{d+1}}\der \xi\leq  C\epsilon
\end{align}
since we chose $d$ sufficiently large such that $\gamma\leq \frac{d\alpha}{d+1}+1$. 



We now bound $(I)$. For $s\leq \tau,$ we only have that \eqref{linear bound for y} holds. Using in addition \eqref{exponential bound for der y}, we have that 
\begin{align*}
    \int_{0}^{\tau(y,v)}\frac{V^{\gamma}|Y|^{d-1}|\partial_{v}Y|} {(1+|Y|^{d})^{2}}\der z\leq   \int_{0}^{\tau}\frac{Cv^{\gamma}|\partial_{v}Y|}{1+|Y|^{d+1}}\der z&\leq \int_{0}^{\tau}\frac{C v^{\gamma+\alpha-1}(e^{\frac{z}{\epsilon}}-1)}{1+\frac{v^{\alpha(d+1)}}{\epsilon^{d+1}}|z-\tau|^{d+1}}\der z\\
    &\leq \int_{0}^{\tau}\frac{C v^{\gamma-1}(e^{\frac{z}{\epsilon}}-1)}{1+\frac{v^{d\alpha}}{\epsilon^{d+1}}|z-\tau|^{d+1}}\der z.
\end{align*}
We now make the change of variables 
\begin{align*}
    \xi&=\frac{v^{\frac{\alpha d}{d+1}}}{\epsilon}(\tau-z);\\
     \der \xi&=-\frac{v^{\frac{\alpha d}{d+1}}}{\epsilon}\der z;\\
     \frac{z}{\epsilon}&=\frac{\tau-\epsilon v^{-\frac{\alpha d}{d+1}}\xi}{\epsilon}\leq \frac{\epsilon \xi +\tau}{\epsilon}= \xi+\frac{\tau}{\epsilon}\leq\xi+1.
\end{align*}
Thus, using that $e^{x}-1\leq Cx,$ when $x\in[0,1]$, we further obtain that
\begin{align}\label{bound i}
    (I)\leq \int_{0}^{\infty}\frac{C\epsilon v^{\gamma-\frac{d\alpha}{d+1}-1}\frac{z}{\epsilon}}{1+|\xi|^{d+1}}\der \xi\leq \int_{0}^{\infty}\frac{C\epsilon v^{\gamma-\frac{d\alpha}{d+1}-1}(\xi+1)}{1+|\xi|^{d+1}}\der \xi\leq C\epsilon,
\end{align}
since $v\geq 1$ and $\gamma\leq \frac{d\alpha}{d+1}+1$. Combining \eqref{bound i} with \eqref{bound ii}, we obtain that \eqref{linear ode part 2} holds.
We denote by 
\begin{align*}
    B(s):=\int_{0}^{s}\frac{L\gamma V^{\gamma-1}}{1+|Y|^{d}}\der z.
\end{align*}
Notice that by \eqref{linear ode part 1}, it holds that $0\leq B(s)\leq CL$, for all $s\in[0,t_{1}]$.
From \eqref{der v v}, we obtain that
\begin{align*}
    \partial_{s}\big(e^{B(s)}\partial_{v}V(s)\big)\geq e^{B(s)}\frac{LdV^{\gamma}|Y|^{d-2}Y\partial_{v}Y}{(1+|Y|^{d})^{2}}\geq -e^{B(s)}\bigg|\frac{LdV^{\gamma}|Y|^{d-2}Y\partial_{v}Y}{(1+|Y|^{d})^{2}}\bigg|\geq -C(L)\bigg|\frac{V^{\gamma}|Y|^{d-2}Y\partial_{v}Y}{(1+|Y|^{d})^{2}}\bigg|,
\end{align*}
for some constant $C(L)>0$ which depends on $L$. Making use of \eqref{linear ode part 2}, we further obtain that
\begin{align*}
e^{B(s)}\partial_{v}V(s)-1\geq -C(L)\epsilon
\end{align*}
and since by \eqref{linear ode part 1} it holds that $0\leq B(s)\leq CL\epsilon$, for all $s\in[0,t_{1}]$, then there exists $\epsilon_{1}\in(0,1)$, which depends on $L$, such that for all $\epsilon\leq \epsilon_{1}$, the following holds
\begin{align*}
    \partial_{v}V(s)\geq e^{-B(s)}(1-C(L)\epsilon)\geq e^{-CL\epsilon}(1-C(L)\epsilon)\geq \frac{3}{4},
\end{align*}
for all $s\in[0,t_{1}]$. We can thus iterate the argument in order to obtain that for all times $t\geq 0$, we have that $\partial_{v}V(y,v,t)\geq \frac{1}{2} \geq 0$ and $\partial_{v}Y(y,v,t)\leq \frac{\alpha}{2}v^{\alpha-1}(1-e^{\frac{t}{\epsilon}})\leq 0$. This concludes our proof.
 \end{proof}

\end{document}
